\documentclass{article}
\usepackage{ctex}
\begin{document}

以下是本手册内容的描述：

第 1 章 — 关于本手册。概述《英特尔® 64 位与 IA-32 架构软件开发人员手册》所有卷册。同时介绍这些手册中的符号约定，并列出程序员和硬件设计人员可能感兴趣的相关英特尔手册与文档。

第 2 章 — 英特尔® 64 位与 IA-32 架构。介绍英特尔 64 位与 IA-32 架构以及基于这些架构的英特尔处理器家族。同时概述这些处理器的共同特性，并简要介绍英特尔 64 位与 IA-32 架构的发展历史。

第 3 章 — 基本执行环境。介绍内存组织模型，并描述应用程序使用的寄存器集。

第 4 章 — 数据类型。描述处理器识别的数据类型和寻址模式；概述实数与浮点数格式以及浮点异常。

第 5 章 — 指令集摘要。列出所有英特尔 64 位与 IA-32 指令，按技术组分类。

第 6 章 — 过程调用、中断与异常。描述过程堆栈以及为过程调用、中断和异常处理提供的机制。

第 7 章 — 通用指令编程。描述对基本数据类型、通用寄存器和段寄存器进行操作的基本加载/存储、程序控制、算术和字符串指令；同时描述在保护模式下执行的系统指令。

第 8 章 — x87 FPU 编程。描述 x87 浮点单元（FPU），包括浮点寄存器和数据类型；概述浮点指令集，并描述处理器的浮点异常条件。

第 9 章 — 英特尔® MMX™ 技术编程。描述英特尔 MMX 技术，包括 MMX 寄存器和数据类型；同时概述 MMX 指令集。

第 10 章 — 英特尔® 流式 SIMD 扩展（英特尔® SSE）编程。描述 SSE 扩展，包括 XMM 寄存器、MXCSR 寄存器和打包单精度浮点数据类型；概述 SSE 指令集，并提供访问 SSE 扩展的编码指南。

第 11 章 — 英特尔® 流式 SIMD 扩展 2（英特尔® SSE2）编程。描述 SSE2 扩展，包括 XMM 寄存器和打包双精度浮点数据类型；概述 SSE2 指令集，并提供访问 SSE2 扩展的编码指南。本章还描述了 SSE 和 SSE2 指令可能生成的 SIMD 浮点异常，并为操作系统和应用程序代码集成 SSE 和 SSE2 扩展支持提供通用指南。

第 12 章 — 英特尔® 流式 SIMD 扩展 3（英特尔® SSE3）、补充流式 SIMD 扩展 3（SSSE3）、英特尔® 流式 SIMD 扩展 4（英特尔® SSE4）和英特尔® AES 新指令（英特尔® AES-NI）编程。概述 SSE3 指令集、补充 SSE3、SSE4、AESNI 指令，并提供访问这些扩展的编码指南。

第 13 章 — 使用 XSAVE 功能集管理状态。描述 XSAVE 功能集指令，并解释软件如何启用 XSAVE 功能集及支持 XSAVE 的功能。

第 14 章 — 英特尔® AVX、FMA 和英特尔® AVX2 编程。概述英特尔® AVX 指令集、FMA 和英特尔® AVX2 扩展，并提供访问这些扩展的编码指南。

第 15 章 — 英特尔® AVX-512 编程。概述英特尔® AVX-512 指令集扩展，并提供访问这些扩展的编码指南。

第 16 章 — 英特尔® AVX10 编程。概述英特尔® AVX10 指令集扩展，并提供访问这些扩展的编码指南。

第 17 章 — 英特尔® 事务同步扩展编程。描述支持锁消除技术的指令扩展，旨在提高具有竞争锁的多线程软件的性能。

第 18 章 — 控制流强制技术。概述控制流强制技术（CET），并提供访问这些扩展的编码指南。

第 19 章 — 英特尔® 高级矩阵扩展编程。概述英特尔® 高级矩阵扩展，并提供访问这些扩展的编码指南。

第 20 章 — 输入/输出。描述处理器的 I/O 机制，包括 I/O 端口寻址、I/O 指令和 I/O 保护机制。

第 21 章 — 处理器识别与功能确定。描述如何确定 CPU 类型以及处理器中可用的功能。

附录 A — EFLAGS 交叉参考。总结 IA-32 指令如何影响 EFLAGS 寄存器中的标志位。

附录 B — EFLAGS 条件码。总结条件跳转、移动和“按条件码设置字节”指令如何使用 EFLAGS 寄存器中的条件码标志（OF、CF、ZF、SF 和 PF）。

附录 C — 浮点异常摘要。总结 x87 FPU 浮点指令和 SSE/SSE2/SSE3 浮点指令引发的异常。

附录 D — 编写 SIMD 浮点异常处理程序指南。提供为 SSE/SSE2/SSE3 浮点指令生成的异常编写异常处理程序的指南。

附录 E — 英特尔® 内存保护扩展。概述英特尔® 内存保护扩展，该功能已被弃用，未来处理器将不再支持。

\end{document}